\chapter*{Phụ Lục}
\addcontentsline{toc}{chapter}{Phụ Lục}
\markboth{Phụ Lục}{Phụ Lục}

\begin{truyen}{Gợi ý phát triển 120 tình huống thành sách hoàn chỉnh}{Appendix for expansion, editing, and classroom use}
\chuhoa{B}{ản} thảo hiện có 12 chương và 120 tình huống. Để phát triển thành một cuốn sách hoàn chỉnh theo đúng nhịp của series, có thể chia tiến độ làm việc thành ba lớp song song.
\textit{(The manuscript now contains 12 chapters and 120 situations. To turn it into a finished volume in the same rhythm as the series, the work can be divided into three parallel layers.)}

\textbf{Lớp 1: Biên tập nội dung.} Mỗi tình huống được mở rộng từ 1 dòng thành 3 phần: bối cảnh, xung đột, và điểm bẻ hướng.
\textit{(Layer 1: content editing. Expand each situation into setting, conflict, and turning point.)}

\textbf{Lớp 2: Đồng bộ song ngữ.} Sau khi có bản tiếng Việt đủ rõ, thêm tiếng Anh ngắn theo kiểu học đường, dễ đọc thành tiếng.
\textit{(Layer 2: bilingual alignment. Once the Vietnamese version is stable, add concise school-friendly English.)}

\textbf{Lớp 3: Ứng dụng lớp học.} Cuối mỗi tình huống có thể thêm câu hỏi thảo luận, từ khóa, hoặc gợi ý hoạt động nhóm.
\textit{(Layer 3: classroom use. Add discussion prompts, keywords, or small group activities.)}
\end{truyen}

\begin{baihoc}[Khung triển khai nhanh]
1. Mỗi chương chọn trước 3 tình huống mạnh nhất để viết hoàn chỉnh.\
2. Giữ cân bằng giữa góc nhìn thầy, trò, phụ huynh và tập thể lớp.\
3. Ưu tiên các tình huống có khả năng dùng ngay trong sinh hoạt lớp hoặc tiết chủ nhiệm.\
4. Khi thêm tiếng Anh, nên giữ câu ngắn, nhịp rõ, tránh văn chương quá tải.

\textit{(1. Fully develop the three strongest situations in each chapter first. 2. Balance teacher, student, parent, and group perspectives. 3. Prioritise situations useful in class meetings. 4. Keep the English short and speakable.)}
\end{baihoc}

\begin{tuhocnhanh}
1. Chương nào có thể trở thành một quyển riêng nếu viết sâu? \textit{(Which chapter could become a standalone volume?)}

2. Tình huống nào cần nhân vật xuyên suốt để mạnh hơn? \textit{(Which situations would benefit from recurring characters?)}

3. Có nên thêm phần hoạt động giáo dục cuối mỗi chương không? \textit{(Should each chapter end with educational activities?)}
\end{tuhocnhanh}

\ngancach